\documentclass[conference]{IEEEtran}
\IEEEoverridecommandlockouts

\usepackage{cite}
\usepackage{amsmath,amssymb,amsfonts}
\usepackage{algorithmic}
\usepackage{graphicx}
\usepackage{textcomp}
\usepackage{xcolor}
\usepackage{booktabs}
\usepackage{tabularx}
\usepackage{pgfplots}
\pgfplotsset{compat=1.18}
\definecolor{plotblue}{RGB}{79,111,215}
\definecolor{plotblueb}{RGB}{107,127,215}
\definecolor{plotgreen}{RGB}{19,122,63}
\definecolor{plotgrid}{RGB}{220,226,235}

% Global adjustments to tighten spacing around tables and figures
\setlength{\textfloatsep}{6pt plus 2pt minus 2pt}
\setlength{\floatsep}{4pt plus 2pt minus 2pt}
\setlength{\intextsep}{4pt plus 2pt minus 2pt}
\setlength{\abovecaptionskip}{3pt plus 1pt minus 1pt}
\setlength{\belowcaptionskip}{3pt plus 1pt minus 1pt}

\def\BibTeX{{\rm B\kern-.05em{\sc i\kern-.025em b}\kern-.08em
    T\kern-.1667em\lower.7ex\hbox{E}\kern-.125emX}}

\begin{document}

\title{Native MXFP4 Dataflow for Persistent Gated DeltaNet Decode on an FPGA\\
}

% \author{\IEEEauthorblockN{Tyler Lee\thanks{All contributions are equal for the first 4 authors}}
% \IEEEauthorblockA{\textit{Electrical and Computer Engineering} \\
% \textit{NYU Tandon School of Engineering}\\
% Brooklyn, NY 11201, USA \\
% tl2966@nyu.edu}
% \and
% \IEEEauthorblockN{Jared Qu}
% \IEEEauthorblockA{\textit{Electrical and Computer Engineering} \\
% \textit{Stony Brook University}\\
% Stony Brook, NY 11794, USA \\
% wolfedqu@gmail.com}
% \and
% \IEEEauthorblockN{Anirvan Kothuri}
% \IEEEauthorblockA{\textit{Electrical and Computer Engineering} \\
% \textit{Stony Brook University}\\
% Stony Brook, NY 11794, USA \\
% anirvanssk@gmail.com}
% \and
% \IEEEauthorblockN{Jason Zheng}
% \IEEEauthorblockA{\textit{Electrical and Computer Engineering} \\
% \textit{Stony Brook University}\\
% Stony Brook, NY 11794, USA \\
% zhengjason219@gmail.com}
% \and
% \IEEEauthorblockN{Jason Blocklove}
% \IEEEauthorblockA{\textit{Electrical and Computer Engineering} \\
% \textit{NYU Tandon School of Engineering}\\
% Brooklyn, NY 11201, USA \\
% jason.blocklove@nyu.edu}
% \and
% \IEEEauthorblockN{Ramesh Karri}
% \IEEEauthorblockA{\textit{Electrical and Computer Engineering} \\
% \textit{NYU Tandon School of Engineering }\\
% Brooklyn, NY 11201, USA \\
% rkarri@nyu.edu}
% }

\author{
\IEEEauthorblockN{Authors anonymized for review}
}

\maketitle

\begin{abstract}
Advances in large language models have caused major bottlenecks due to the substantial computation and storage needs required to process long strings of input data as well as the storage needed for all internal key value representations. The key-value (KV) cache for a traditional Transformer model scales linearly with sequence length. Hybrid methods were proposed that combine both linear and full-attention layers within one neural network architecture. However, even with these hybrid architectures, the use of either high precision floating point numbers, such as 32-bit single-precision floating-point (FP32), or standard low precision floating point numbers (e.g., FP8) quickly exhausts available on-chip memory. We propose a novel architecture that uses a native 4-bit Open Compute Project (OCP) Microscaling (MXFP4) datapath that reduces the amount of space required to store fixed-size linear state matrices. Testing of our proposed MXFP4 architecture was done on an AMD Alveo U55C FPGA field-programmable gate array (FPGA). The results from our evaluation indicated that the MXFP4-based implementation achieved a 5.537$\times$ speedup and a 216.011$\times$ improvement in energy efficiency relative to the H100 baseline, as well as an 86.719\% reduction in recurrent-state storage relative to FP32.
\end{abstract}


\begin{IEEEkeywords}
Large Language Models, Hardware Accelerator, FPGA, MXFP4, Linear Attention, Hybrid Transformers, Memory Hierarchy Partitioning, Alveo U55C.
\end{IEEEkeywords}

\section{Introduction}
Autoregressive decoding in large language models introduces measurable latency at inference time, as each generated token requires sequential access to model state. To understand the reason for this kind of delay, we must first realize the framework being used at the front of modern LLM operations. Many of the most popular models, such as ChatGPT and Claude, were founded through the transformer model architecture~\cite{b2}, a completely software-based mathematical technique that processed data sequences simultaneously using a ``Self-Attention'' mechanism that decides how much one word would relate to another to form context and reference. When a new word is generated, the model retrieves the previous words, or ``tokens''. This bottleneck is mitigated by storing historical token data in a structured mathematical format into external memory, which would be called a key-value cache (KV cache)~\cite{b6}.

The present industry’s trend of utilizing a hybrid layer structure does address the exploding usage of having only full attention layers which mitigates an aspect of the quadratic scaling nature of those layers. However, that structure doesn’t eliminate the nature of the growing complexity, but rather reduces the growth of complexity as compared to a scenario of pure full attention layers~\cite{b9}. Furthermore, traditional implementations utilizing high-precision data formats such as FP32 will quickly saturate the limited memory on these kinds of field-programmable gate arrays (FPGA)~\cite{b10}, an integrated circuit that hosts configurable hardware blocks that allow the user to make their own custom digital logic circuits. When these state matrices start to exceed the local capacity, the hardware must go on to use external bandwidth memory.

To break this memory wall and improve efficiency of current integer standards, we implement a MXFP4 datapath, a 4-bit microscaling floating point data format that reduces model size by four times by sharing a 8-bit scaling factor, onto a Alveo U55C accelerator to create an efficient connection between the local memory and external memory. These specific memory specifications have been studied and observed thoroughly through researchers such as C. Wolters, X. Yang, U. Schlichtmann, and T. Suzumura~\cite{b5}. Bypassing the memory wall is achievable by treating the URAM and BRAM as the permanent execution area for those compressed $S$ matrices and reserving the high bandwidth memory (HBM) for the growing KV caches of the full attention layers as explored by other research papers~\cite{b10}.

By maximizing the efficiency of the data format itself, we can better utilize the limited on-chip storage available. Numerical states can be packed into extremely compact 4-bit Open Computer Project (OCP) microscaling formats through the specific utilization of 4-bit E2M1 elements, a format representation that distributes its four bits into a sign bit, exponent bits, and a precision bit, that are bounded by a shared 8-bit E8M0 block scale, a 8-bit format representation consisting of eight exponent bits with zero precision mantissa bits~\cite{b12}. The processing elements are able to perform continuous linear attention updates without having to request a time-consuming demand to the external memory bus. By focusing on one piece of a hybrid attention model, we can improve the performance and efficiency of the entire implementation that is being used. In a hybrid model like the previously mentioned 3:1 structure, we can look at the primary building block that makes up 75\% of the core pipeline and focus entirely on improving what makes up the majority of the model which would naturally optimize the entire system as a whole~\cite{b9}.

The primary objective of this research is to maximize the use of low-latency on-chip memory resources that maintains the precision needed for long context models. By mapping 4-bit MXFP4 multipliers directly onto the lookup table (LUT) resources of the FPGAs, parallel computation is achieved~\cite{b12}. This paper effectively introduces a new data packaging optimization method that can break through the modern memory wall and unlock higher potential of hybrid linear attention architectures on high-end FPGA fabric~\cite{b5}.

\section{Related Work}

\subsection{Memory Bottlenecks in LLM Decode}
Transformers are typically limited by memory movement in addition to compute throughput. The growth of the KV cache with increasing sequence length reduces both the ability to use high-bandwidth memory resources efficiently and the overall accelerator utilization at longer context lengths. Systems such as vLLM improve LLM serving throughput by managing KV-cache memory more efficiently \cite{b4}, while other studies of LLM serving have identified memory bandwidth and cache movement as primary limitations \cite{b1,b7}.

\subsection{Linear Attention and Gated DeltaNet}
As sequence length increases, the KV cache grows proportionally, consuming additional memory bandwidth and storage resources~\cite{b7}, thus driving the development of another mathematical alternative called linear attention, like the Gated DeltaNet (GDN) model~\cite{b3}. This linear-time sequence model compresses the history of data down into a single fixed grid known as a recurrent state matrix $S$~\cite{b3}. The tradeoff with the linear-attention layer is that historical token details are heavily compressed. Nevertheless, by combining both approaches and stacking them, researchers have made the token go through linear layers first then into a full-attention layer, known as the 3:1 structure block, to more efficiently process previous history while maintaining data integrity, becoming the present industry’s forced solution.

Linear and hybrid attention architectures reduce reliance on large KV caches during long-context decoding by replacing some attention layers with fixed-state recurrent layers. The per-token complexity of full attention is $\mathcal{O}(N)$ with respect to sequence length $N$, which causes total complexity over an entire sequence to scale quadratically as $\mathcal{O}(N^2)$. Linear attention reduces per-token complexity to $\mathcal{O}(1)$, so total sequence complexity scales as $\mathcal{O}(N)$, lowering the number of multiply-accumulate (MAC) operations needed for long token histories~\cite{b3}. Hybrid architectures such as the 3:1 structure --- which routes tokens through three linear-attention layers before each full-attention layer --- do not eliminate quadratic growth entirely, but reduce it relative to pure full-attention designs~\cite{b9}. One example of a linear-attention architecture is the Gated DeltaNet family of models, which compresses token history into a fixed-size recurrent state matrix rather than storing each previous key and value vector explicitly~\cite{b3}. A fixed-size recurrent state provides a particular advantage for FPGA accelerators: if the state remains resident on chip, the design can avoid repeatedly transferring it between on-chip memory and external memory~\cite{b10}.

\subsection{FPGA Acceleration for Attention Models}
This work uses FPGAs as a basis for implementing application-specific architectures, specifically in terms of dataflow architecture. These devices provide flexibility to map memory banks based upon the specific requirements of the workload or application and enable pipelined operation and custom mapping of arithmetic units.

Prior work on using FPGAs for persistent-state architectures demonstrated that storing linear attention states on-chip improves memory efficiency when performing decode-related workloads \cite{b10}. Building upon this approach, we focus on a native MXFP4 datapath that directly targets low-precision persistent-state computation rather than relying on a BF16-style high-precision datapath.


\subsection{Low-Precision and Microscaling Formats}
Reduced-precision data formats reduce memory usage and arithmetic cost. However, a fixed low-bit-width representation may limit the ability to represent large ranges in individual elements. To improve upon fixed-format representations, microscaling provides a method of using a single exponent value that is common to multiple elements \cite{b12}. MXFP4 utilizes 4-bit E2M1 values with a common scale applied to a small group of values, or block \cite{b13}. This results in compact storage for each individual element and allows the entire block to store a larger local dynamic range \cite{b14}. Due to these characteristics, MXFP4 represents a strong candidate for hardware-based implementations on FPGAs.

The E2M1 element multiply in MXFP4 can be directly implemented in FPGA LUT logic, while the block exponent determines the alignment before accumulation. As such, MXFP4 has been selected as the primary datapath and is used as the basis for comparison in this work, while INT4 is used as the primary alternative comparison for evaluation purposes \cite{b15}.

\section{Problem Formulation and Workload Definition}

The primary focus of this paper is a batch-size-one Gated DeltaNet decode layer during autoregressive inference. This focus isolates the recurrent-state update, which is the portion of a hybrid linear-attention model most directly influenced by persistent on-chip storage. In recurrent linear attention, sequence history is stored in a fixed-size state rather than in a KV cache that grows with sequence length. Because of this property, Gated DeltaNet is naturally suited to an FPGA dataflow in which the recurrent state can remain in on-chip memory across successive decode steps~\cite{b3,b10}.

The design motivation comes from the latency-sensitive nature of autoregressive decode. Each newly generated token depends on the prior state produced by the previous decoded token. Therefore, multiple tokens cannot be processed concurrently across decode steps. As a result, repeatedly reading and writing model state can become a dominant performance cost. For batch-size-one Gated DeltaNet decode, moving the recurrent state through HBM at each token can create a memory-bound bottleneck rather than a purely compute-bound one~\cite{b10}.

This motivates a persistent-state datapath rather than a conventional off-chip state-transfer datapath. The architecture is designed so that the recurrent state remains in on-chip FPGA memory while the token-level computation is performed around it. This organization reduces repeated HBM movement, exposes parallelism within each single-token update, and allows the arithmetic path to be specialized for the recurrent-state operation. The use of MXFP4 further supports this architecture by reducing the storage required for the recurrent state while allowing the multiply path to be mapped efficiently onto FPGA logic.

Accordingly, this study examines a single Gated DeltaNet layer rather than full-model LLM inference. A full-model formulation would include embeddings, feed-forward layers, normalization, attention layers, and host-device scheduling overheads. These components are important for deployment, but including them would obscure the specific memory-traffic and arithmetic behavior of the recurrent update kernel. Restricting the workload to a single recurrent decode layer makes it easier to attribute latency, energy, and storage effects to the persistent-state MXFP4 datapath.

Alternative workload definitions and architectures were also considered. A prefill-oriented workload would emphasize parallel processing across an input sequence, whereas this work targets the token-by-token dependence that defines autoregressive decode. A full hybrid-attention workload would include both recurrent and attention layers, but that scope would address broader model-level architecture questions rather than the recurrent-state storage problem isolated here~\cite{b3,b8,b9,b11}. A conventional GPU or FPGA design that repeatedly moves the recurrent state through HBM would preserve a memory-bound execution pattern. A BF16 or FP8 datapath would reduce some numerical risk, but it would require substantially more on-chip storage than MXFP4. An INT4 datapath was also considered as a lower-precision alternative, but MXFP4 better matches the goal of combining compact storage with block-scaled floating-point dynamic range.

More generally, this problem formulation applies to sub-quadratic sequence models that use fixed-size state spaces rather than sequence-length-dependent Transformer KV caches. This structure reduces the growth of stored sequence history, but it also exposes a recurrent-state movement bottleneck during decode that can be addressed with specialized hardware~\cite{b2,b10}.

\section{Hardware Architecture}

Figure~\ref{fig:mxfp4_gdn_datapath} summarizes the proposed native MXFP4 persistent-state Gated DeltaNet decode datapath. The accelerator receives a query vector, key vector, value vector, gate, and update parameter for each generated token. These inputs are used to read the prior recurrent state, compute the gated delta-rule update, write the updated state back to on-chip memory, and generate the output vector.

The datapath is organized around persistent storage of the recurrent state in on-chip FPGA memory. This organization allows the design to exploit spatial parallelism within a single-token update while reducing unnecessary off-chip movement of the recurrent state. The native MXFP4 arithmetic path is reused during prediction, state update, and output generation. The default configuration uses 32 attention heads with a head dimension of 128, matching the target Gated DeltaNet layer used throughout the evaluation.

\begin{figure}[t]
\centering
\includegraphics[width=\linewidth]{mxfp4_gdn_datapath_paper_fit.png}
\caption{Overview of the proposed native MXFP4 persistent-state Gated DeltaNet decode datapath. Token inputs are packed into MXFP4 blocks, the recurrent state remains resident in on-chip BRAM/URAM, and native E2M1 LUT multiplication with block-scale alignment and fixed-point accumulation is reused across prediction, state update, and output generation.}
\label{fig:mxfp4_gdn_datapath}
\end{figure}

\subsection{Persistent Recurrent State}
As opposed to an existing KV cache that expands in size as sequence length increases, the recurrent state of Gated DeltaNet consists of a fixed-size recurrent-state matrix \cite{b3}. The recurrent-state matrix stores the compressed information of previous tokens, functioning as a compact associative memory that addresses the standard LLM memory wall. The recurrent-state is updated after every decode cycle.

The recurrent state $S_t$ is stored within partitions of on-chip FPGA memory banks, allowing access to different parts of the state matrix in parallel. As a result, on-chip storage of the recurrent state reduces latency and energy consumption since each decode step needs to load the prior state, modify it, and save the modified version. However, accessing the full state matrix repeatedly off chip and back on chip would become a resource intensive process that would worsening the memory-bound limits of standard decode paths.

The recurrent state is the primary persistent storage entity in the Gated DeltaNet decode layer. By using MXFP4, the amount of storage required for each element is reduced to 4 bits per element plus a small amount of shared block-scale metadata. In comparison, BF16 requires 16 bits per element, while FP32 requires 32 bits per element. Consequently, MXFP4 reduces the storage requirements for recurrent-state elements by nearly 4 times compared with BF16 and 8 times compared with FP32, before accounting for the small overhead associated with block-scale metadata, as shown in Table~\ref{tab:state_traffic_results}.

\begin{table}[t]
\centering
\caption{Latency comparison for batch-one Gated DeltaNet decode.}
\label{tab:latency_results}
\footnotesize
\renewcommand{\arraystretch}{1.0}
\setlength{\tabcolsep}{6pt}
\begin{tabularx}{\linewidth}{Xcc}
\toprule
Design & Latency ($\mu$s/token) & Speedup vs. H100 \\
\midrule
H100 baseline & 285.0 & 1.00$\times$ \\
FPGA baseline & 63.2 & 4.51$\times$ \\
Proposed MXFP4 & 51.472 & 5.54$\times$ \\
\bottomrule
\end{tabularx}
\end{table}

\subsection{Gated DeltaNet Update Path}
Every decode step uses the current key and value to update the recurrent-state. First, the prior state is multiplied by the key vector to create a prediction based on delta-rule dynamics \cite{b3}:
\begin{equation}
r_t = S_{t-1} k_t
\end{equation}

The prediction is what the prior state already holds for the upcoming key. The hardware performs a comparison between the prediction and the new value vector:
\begin{equation}
e_t = r_t - v_t
\end{equation}

The error ($e_t$) is created due to a difference between what was stored by the prior state and what information is presented by the new token. This error is scaled by the beta input and is used to update the state:
\begin{equation}
S_t = S_{t-1} - \beta_t e_t k_t^T
\end{equation}

This state update is a rank-one update. It changes the state by applying both the direction of the key and the error against the value.

Following the recurrent state update, the design utilizes the query vector and the recurrent state to produce the intermediate output vector:
\begin{equation}
o_t = S_t q_t
\end{equation}

As this is a Gated DeltaNet layer, the design applies an output gate to the intermediate output. The gated output can be represented as:
\begin{equation}
y_t = g_t \odot o_t
\end{equation}
where $g_t$ is the output gate and $\odot$ is elementwise multiplication. The output gate controls how much of the recurrent-state output is forwarded to subsequent layers. As such, the gating mechanism provides a means of regulating the recurrent-state contribution to each generated token.

\subsection{Five-Step Token Pipeline}
Gated DeltaNet could naively be implemented over three steps. First, an old-state $S_{t-1}$ would have been loaded so that a prediction could be made regarding how to retrieve data using a key vector. Next, the retrieved data would have been used to correct/update the state. Lastly, an updated-state would have been loaded so that an output could be produced using a query vector. The resulting design would be inefficient for an FPGA implementation because the largest structure (the state-matrix) must be traversed completely by every token.

However, we were able to optimize the sequence of calculations such that only partial traversal of the full state-matrix occurs, implementing a fused dataflow framework tailored for linear attention bottlenecks \cite{b10}. We did this by limiting use of reusable terms derived from state-data used in computing retrieval terms to those computations used in generating an output; therefore, no additional full read of an updated-state needs occur for output generation purposes. Every token goes through five states.

\begin{table}[t]
\centering
\caption{Five-step token pipeline in the proposed datapath.}
\label{tab:five_step_token_pipeline}
\scriptsize
\begin{tabular}{clp{0.56\linewidth}}
\toprule
Step & Module & Function \\
\midrule
1 & Input preparation & Put each token in its corresponding scaled MXFP4-block. \\
2 & State-read & Load current-state $S_{t-1}$ from parallel memory-banks to calculate retrieval terms for the delta-rule. \\
3 & Delta-update & Calculate beta-scaled error and state-corrective terms. \\
4 & Write-back & Write updated-state $S_t$ directly back to the memory-banks where they reside. \\
5 & Output-generation & Generate final-token-output $o_t$ by adding query-vectors with previously-generated-state-derived-terms. \\
\bottomrule
\end{tabular}
\end{table}

First, the input preparation module puts each token in its corresponding scaled MXFP4-block. Second, the state-read module loads up the current-state $S_{t-1}$ from parallel memory-banks in order to calculate retrieval terms for the delta-rule. Third, the delta-update module calculates beta-scaled error and state-corrective terms. Fourth, the write-back phase writes updated-state $S_t$ directly back to memory-banks where they reside. Fifth, the output-generation module generates final-token-output $o_t$ by adding query-vectors with previously-generated-state-derived-terms. As such, we eliminate a second-full read of-state.

Overall, our optimization reduces total-state-traversal work from approximately three full passes to less than two by collapsing these passes into a tighter, five-phase pipelined datapath design \cite{b10}.

\subsection{MXFP4 Quantization Method}

The proposed design utilizes MXFP4 as the primary arithmetic format. MXFP4 is a representation where each value is represented as a 4-bit E2M1 element (one sign bit, two exponent bits and one mantissa bit), along with a common E8M0 block exponent. Unlike the other formats which assign an independent exponent to each value, the format provides a single exponent for all values within a block. Assigning one exponent to all values within a block provides a reduction in memory usage while providing greater dynamic range than would be possible with a flat INT4 format.

Each value in a block is approximated as the product of the shared block scale, and a quantized E2M1 element. The shared exponent captures the scale of the entire block, while the E2M1 element captures the signed mantissa and exponent information for each individual element.

The hardware datapath defaults to an MXFP4 block size of 32. We validated the recurrent-state quantization logic with deterministic synthetic test vectors generated in Python, and verified bit-exact agreement with the MXFP4 hardware implementation. This verifies functional correctness, but synthetic data does not capture the extreme-magnitude outliers that are common in real LLM activations. Hence our conclusions are limited to numerical fidelity and hardware correctness. To evaluate the effect on model level metrics, e.g., perplexity, would require live activation traces from models like Qwen3-Next, and is out of scope of this work.

\begin{table}[t]
\centering
\caption{Recurrent-state storage and off-chip traffic accounting for MXFP4 $B=32$.}
\label{tab:state_traffic_results}
\footnotesize
\renewcommand{\arraystretch}{1.0}
\setlength{\tabcolsep}{6pt}
\begin{tabularx}{\linewidth}{Xr}
\toprule
Metric & Value \\
\midrule
MXFP4 $B=32$ recurrent-state storage & 278528 bytes \\
Reduction vs FP32 & 86.719\% \\
Reduction vs BF16 & 73.438\% \\
Naive three-pass off-chip state traffic & 835584 bytes/token \\
Persistent steady-state off-chip state traffic & 0 bytes/token \\
\bottomrule
\end{tabularx}
\end{table}

\begin{table}[t]
\centering
\caption{Arithmetic formats used in the proposed datapath.}
\label{tab:data_formats}
\footnotesize
\renewcommand{\arraystretch}{1.0}
\setlength{\tabcolsep}{6pt}
\begin{tabularx}{\linewidth}{llX}
\toprule
Tensor / Value & Format & Notes \\
\midrule
Weights & MXFP4 & Block size 32 \\
Activations & MXFP4 & Dynamic per token \\
Recurrent state & MXFP4 & Stored on chip \\
Beta & INT8 & Scaling parameter \\
Partial products & Q4.3 & From E2M1 multiply \\
Accumulator & INT24 & Block-aligned accumulation \\
Final output & FP16 & Output format \\
\bottomrule
\end{tabularx}
\end{table}

\subsection{FP4 Element Multiply}
The FP4 element multiply is implemented as a native E2M1 operation, rather than by converting operands into BF16 or FP32. The multiply path consists of: sign XOR on the sign bits, exponent addition with bias adjustment, mantissa multiplication, normalization; the normalized result of the element multiply is then aligned based upon the shared block exponents, before being accumulated. This LUT-based implementation reduces the physical area of the inner multiplier path and avoids the need for large floating-point DSP blocks for the four-bit element multiply.

MXFP4 quantizes elements differently than flat INT4. Flat INT4 uses extremely compact values, however it also lacks local dynamic-range flexibility. MXFP4 maintains the same 4-bit element width as INT4 while utilizing the block exponents to adapt to local tensor magnitude. The synthetic numerical results in this paper demonstrate that MXFP4 retains better output directionality than the currently used INT4 fallback under the current deterministic test conditions.

\subsection{LUT-Based E2M1 Multiplication}
LUT-based E2M1 multipliers are utilized in the FPGA LUT logic as they avoid conversion of operands from BF16 (or FP32) format and provide a direct path to E2M1 multiplications. Each E2M1 is split into three parts: Sign, Exponent, and Mantissa. An XOR operation is performed on the Sign field. A biased sum of Exponents is calculated. The Mantissas are multiplied together and then normalized. The resulting fixed point partial-product is prepared for accumulation. Additionally, this process minimizes the physical area of the inner Multiplier Block and eliminates the requirement for large floating-point DSP Blocks that would be necessary to process four-bit inputs.

\subsection{Block-Exponent Alignment}
Since MXFP4 employs shared block-scales, the magnitudes of the products of two blocks can vary. After generation of partial-products, these values are therefore adjusted with regard to their block-exponents prior to summation. If two blocks have exponents $s_a$ and $s_b$, then
\begin{equation}
s_{prod} = s_a + s_b
\end{equation}
is the combined exponent of the product. Next, the partial-product will be shifted-left so that it is placed in the correct range for summation with an accumulator. This approach ensures the numeric correctness of the scaled values in MXFP4.Three design parameters determine the behavior of the overall design architecture: Block Size ($B$), Key Parallelism ($P_K$) and Value Parallelism ($P_V$).

\section{Experimental Methodology}

\subsection{Experimental Setup}
The proposed accelerator is evaluated as a simulation-only FPGA design targeting the AMD Alveo U55C platform. The target device is the \texttt{xcu55c-fsvh2892-2L-e}, and the design is synthesized using Vitis HLS with a target clock period of 4 ns, corresponding to 250 MHz. Post-synthesis and post-implementation reports are generated through the AMD FPGA tool flow to estimate timing, resource utilization, and power.

The evaluated kernel implements one Gated DeltaNet decode layer for batch-one autoregressive inference. The kernel receives one token's query, key, value, beta, and gate inputs, updates the recurrent state, and produces one output vector. Projection layers, tokenizer logic, full-model scheduling, and multi-layer execution are outside the scope of this evaluation \cite{b3,b10}.

\subsection{Reference Models and AI-Assisted Verification}
The software reference is implemented in Python for high-precision matrix validation. The FP32 model provides the baseline for Gated DeltaNet recurrence, and the custom MXFP4 reference defines the quantized behavior expected from the RTL design.

Deterministic test vectors are generated using a fixed seed to ensure that the simulation results can be reproduced. These vectors include token inputs, recurrent state values, quantized MXFP4 operands, and expected output values. Generative AI assistance was used to draft standard hardware verification testbenches, data-formatting scripts, and baseline structural code. The generated code was reviewed, modified, and verified for functional correctness by the authors.

\subsection{Quantization Configuration}
The evaluated design uses MXFP4 as the primary low-precision format. Each element is represented using a four-bit E2M1 value, and each block shares an E8M0 scale with a default block size of $B=32$ \cite{b12}. Quantization uses deterministic block-scaled conversion and round-to-nearest-even behavior so that the software reference and HLS implementation can be compared bit-exactly.

Weights use MXFP4 with block size 32, activations use dynamic per-token MXFP4 with block size 32, and the recurrent state is stored in MXFP4 on chip. The beta value is represented using symmetric INT8. Element multiplication produces Q4.3 fixed-point partial products, block accumulation uses INT24, and the final pre-dequantized output is represented in FP16, Table~\ref{tab:data_formats}.

\subsection{Hardware Configuration and LUT-Based Datapath}
The default hardware configuration uses block size $B=32$, key-direction parallelism $P_K=16$, and value-direction parallelism $P_V=8$.

The multiplier performs sign XOR, exponent addition with bias adjustment, mantissa multiplication, normalization, and subnormal handling before producing a Q4.3 fixed-point partial product. Block-exponent alignment then shifts and accumulates these partial products using the shared MXFP4 block scales.

\subsection{Simulation and Synthesis Flow}
The evaluation uses a hierarchical verification flow based on the Vitis HLS toolchain. First, HLS C simulation verifies functional correctness against the quantized MXFP4 software reference. Second, HLS C synthesis estimates latency, initiation interval, and initial resource mapping. The synthesis reports are analyzed to confirm that the native E2M1 multiply primitive maps to LUT logic rather than DSP blocks.

The HLS implementation makes use of explicit synthesis pragmas to direct the synthesized hardware. Pipelining is performed on the token-loading, prediction, state-update, output-storage, and block-exponent accumulation loops with a target loop initiation interval of $\mathrm{II} = 1$. Loop unrolling and tiled loop structure are both employed to expose the key- and value-parallel lanes managed by $P_K$ and $P_V$, respectively. In addition to being accessed in parallel, the recurrent-state buffer is partitioned along the key and value dimensions, and it is explicitly mapped onto dual-port BRAM. The top-level function uses AXI memory interfaces with separate memory bundles for the major tensor streams.

Third, RTL cosimulation validates how the RTL behaves across a number of autoregressive decode steps, as such confirming that recurrent updates to its persistent recurrent state occur with no need for recurring off-chip recurrent-state traffic. Fourth, post-implementation Vivado reports provide the authoritative estimates of timing closure, total power consumption, and resource utilization; the final utilization metrics report LUTs, BRAM, URAM where applicable, and DSP usage.

\subsection{Baselines}
The proposed design is evaluated against two primary baselines to contextualize latency and resource efficiency. The first baseline is an NVIDIA H100 PCIe GPU using FP32/BF16 arithmetic as a general-purpose accelerator reference for autoregressive LLM inference. The second baseline is a prior FPGA persistent-state Gated DeltaNet accelerator using BF16 arithmetic and standard DSP mapping \cite{b10}. This baseline is critical for architectural comparison because it targets the same linear-attention decode bottleneck but does not use the native MXFP4 datapath.

Speedup is computed by dividing the baseline latency by the latency of the proposed design:
\begin{equation}
\text{Speedup} = \frac{T_{\text{baseline}}}{T_{\text{ours}}}
\end{equation}

Energy per token is computed to demonstrate the power efficiency of the four-bit datapath:
\begin{equation}
E_{\text{token}} = P \times T_{\text{token}}
\end{equation}
where $P$ is the post-implementation power in watts and $T_{\text{token}}$ is the absolute token latency in seconds.

\subsection{Evaluation Metrics}
The architectural evaluation is divided into three metric categories: performance, hardware utilization, and numerical fidelity.

Performance is quantified primarily by latency per decoded token, measured in microseconds ($\mu s$), and the corresponding energy per token, derived from post-implementation power estimates in watts ($W$). Hardware utilization metrics focus on the absolute consumption and percentage of available resources on the Alveo U55C. LUT utilization is tracked to validate the area efficiency of the native FP4 datapath, while BRAM utilization is tracked to validate the storage feasibility of the persistent MXFP4 recurrent state. DSP usage is monitored to distinguish the LUT-based FP4 multiply primitive from other wider fixed-point arithmetic that may still be inferred by HLS.

Numerical fidelity is evaluated using output cosine similarity and state relative L2 error. These synthetic metrics compare the MXFP4 output against the high-precision FP32 reference model and quantify the isolated error introduced by the four-bit MXFP4 quantization pipeline.

\subsection{Provenance and Reproducibility}
To ensure scientific rigor and reduce transcription errors, reported performance metrics are extracted automatically from generated toolchain reports rather than entered manually. Latency values are extracted from HLS and RTL cosimulation logs, while resource utilization and power estimates are parsed directly from Vivado post-implementation reports.

These extracted values are aggregated into a machine-readable JSON data file and mapped to their source reports through a tracked provenance file. This automated extraction flow ensures that every table and figure in the manuscript can be traced back to a source file, line number, extraction script, timestamp, and Git revision.

\section{Results}



\begin{table}[t]
\centering
\caption{Energy per token comparison for the cited baselines and the proposed MXFP4 design.}
\label{tab:energy_results}
\footnotesize
\renewcommand{\arraystretch}{1.0}
\setlength{\tabcolsep}{6pt}
\begin{tabularx}{\linewidth}{Xcc}
\toprule
Design & Energy (mJ/token) & Efficiency vs. H100 \\
\midrule
H100 baseline & 99.8 & 1.00$\times$ \\
FPGA baseline & 9.5 & 10.51$\times$ \\
Proposed MXFP4 & 0.462 & 216.01$\times$ \\
\bottomrule
\end{tabularx}
\end{table}

\subsection{Latency and Energy}

These results regarding latency and energy support the primary architectural argument made in this document \cite{b5,b10}. Therefore, when comparing both baselines to the proposed MXFP4 design, there were significant reductions in both latency and energy consumption due to maintaining the recurrent state on chip and utilizing a compact native MXFP4 datapath for the decode computations, as shown in Tables~\ref{tab:latency_results} and~\ref{tab:energy_results} \cite{b10,b12}.


\subsection{Timing and Resource Utilization}
The implementation results in Table~\ref{tab:resource_results} show that the designed datapath achieves the 250 MHz timing target with sufficient timing margin on the Alveo U55C. The low LUT and BRAM utilization implies that the FPGA resources are sufficient to support the examined single-layer MXFP4 recurrent workload with native FP4 arithmetic. The reported DSP usage should be viewed as top-level/non-MAC logic rather than evidence that the E2M1 multipliers utilize DSP blocks.

\begin{table}[t]
\centering
\caption{Post-implementation timing and resource utilization on the Alveo U55C.}
\label{tab:resource_results}
\footnotesize
\renewcommand{\arraystretch}{1.0}
\setlength{\tabcolsep}{6pt}
\begin{tabularx}{\linewidth}{Xc}
\toprule
Metric & Result \\
\midrule
HLS estimated Fmax & 342.47 MHz \\
Post-implementation slack & 0.065 ns \\
LUT utilization & 6.64\% \\
DSP utilization & 7.27\% \\
BRAM utilization & 3.67\% \\
\bottomrule
\end{tabularx}
\end{table}

\subsection{Synthetic Numerical Fidelity}
Current numerical-quality results must be viewed as only being used as evidence. Given that realistic Qwen3-Next activation capture is still blocked by dependencies, this version of the manuscript cannot claim final Qwen3-Next perplexity or overall model quality \cite{b3}. As shown in Table~\ref{tab:accuracy_results}, deterministic synthetic test harness, MXFP4 performs better than the INT4 fallback in preserving the direction of outputs. Specifically, INT4 reached an output cosine similarity of 0.967, whereas MXFP4 reached an output cosine similarity of 0.982 for block size = 32 and an output cosine similarity of 0.982 for block size = 16 \cite{b10,b12}.

\begin{table}[t]
\centering
\caption{Synthetic output cosine similarity for INT4 and MXFP4 configurations.}
\label{tab:accuracy_results}
\footnotesize
\renewcommand{\arraystretch}{1.0}
\setlength{\tabcolsep}{6pt}
\begin{tabularx}{\linewidth}{Xc}
\toprule
Configuration & Output cosine similarity \\
\midrule
INT4 fallback & 0.966696 \\
MXFP4, $B=32$ & 0.981823 \\
MXFP4, $B=16$ & 0.982292 \\
\bottomrule
\end{tabularx}
\end{table}


\begin{table}[t]
\centering
\caption{Synthetic recurrent-state stress evidence.}
\label{tab:state_stress_results}
\footnotesize
\renewcommand{\arraystretch}{1.0}
\setlength{\tabcolsep}{6pt}
\begin{tabularx}{\linewidth}{Xcc}
\toprule
Synthetic test & MXFP4 state & MXFP8 state \\
\midrule
1024-token final output cosine & 0.643 & 0.933 \\
1024-token final state rel. L2 & 0.865 & 0.344 \\
256-token boundary output cosine & 0.747 & 0.945 \\
Combined-stress worst output cosine & 0.928 & 0.935 \\
Outlier-stress worst output cosine & 0.968 & 0.973 \\
\bottomrule
\end{tabularx}
\end{table}

\subsection{Recurrent-State Stress}
In addition, Table~\ref{tab:state_stress_results} reports recurrent-state stress data that clarifies the limitations of the current synthetic data. For a synthetic stress test on 1024 tokens, MXFP4 recurrent state produced an output cosine similarity of 0.643 and a final state relative L2 error of 0.865. In comparison, MXFP8 recurrent state produced an output cosine similarity of 0.933 and a final state relative L2 error of 0.344 \cite{b10,b12}. The two results demonstrate that MXFP4 is always the native four-bit arithmetic path, but the stress test also shows that storing the recurrent state in MXFP4 can accumulate state drift over longer decode sequences. MXFP8 recurrent state will always be a safer fallback when accumulation of recurrent-state drift becomes the limiting factor in terms of accuracy, since it trades some storage efficiency for improved numerical stability~\cite{b3,b10,b12}.



\subsection{Storage and Traffic Accounting}
The storage and traffic accounting support the persistent-state data flow claim \cite{b5,b10}. The MXFP4 block-size-32 recurrent state used 278,528 bytes of memory; this represents a total 86.719\% reduction relative to FP32 and 73.438\% reduction relative to BF16 \cite{b10,b12}. A naive three-pass off-chip recurrent-state path would have transferred approximately 835,584 bytes per token \cite{b7,b10}. As such, the persistent-state datapath has zero bytes per token of steady-state off-chip recurrent-state traffic \cite{b7,b10}.



\subsection{Summary of Supported Claims}
Ultimately, these results provide strong evidence supporting both aspects of the claim for a single-layer Gated DeltaNet decode accelerator with simulation-based backing—i.e., hardware and functional behavior \cite{b3,b10}. The most strongly supported claims are RTL cosimulation latency, Vivado timing closure, post-implementation resource usage, post-implementation power, reduced state storage, avoided traffic, and synthetic numerical fidelity \cite{b3,b5,b10,b12}. Claims related to overall model quality should remain pending until we obtain realistic Qwen3-Next activation data \cite{b3,b6,b8,b9,b11}.

\section{Discussion and Limitations}

\subsection{Scalability and BRAM Capacity Ceiling}
The current existing architecture has been created to be optimal for processing tasks in the single-layer single-token inference batch size one configuration. This approach can help maximize on-chip data locality and reduce repeated recurrent-state movement through external memory. Although this configuration works for the evaluated design point, the main scaling restriction is connected with the limited on-chip memory capacity of the Alveo U55C FPGA. For larger recurrent state tensors, larger batch sizes, or expanded model configurations, the amount of memory required by the algorithm may become greater than the on-chip SRAM capacity, making off-chip HBM use necessary. Consequently, DRAM bandwidth limitations may appear again, a bottleneck extensively documented in large-batch LLM serving systems \cite{b1,b4,b7}. The memory-bound nature of autoregressive decode, in which each generated token requires access to model state, makes bandwidth a dominant constraint at increasing batch sizes \cite{b5}. Further research may evaluate tiling and buffer hierarchies across BRAM, URAM, and HBM, a direction also identified in persistent-state FPGA dataflow designs for linear attention \cite{b10}.

\subsection{LUT Routing Congestion}
The LUT-based E2M1 multiply primitive is compact in area, but a full neural layer can instantiate many parallel multiply paths. The evaluated parallelism point uses $P_K=16$ and $P_V=8$, which creates many routes from 6-input lookup tables into the surrounding accumulation and alignment logic. Such logic architecture can place pressure on FPGA interconnect as the design scales, increasing wire length and potentially reducing timing margin. Possible ways to address this problem include floorplanning related logic into pblocks, retiming, and reevaluating the parallelism ratio relative to throughput \cite{b10}.

\subsection{The ``Stubborn Compiler'' Threat (HLS DSP Inference)}
The MXFP4 datapath code is written in C++ and implemented through Vitis HLS. The intended architectural point is that the FP4 element multiply primitive is implemented using lookup tables rather than DSP48E2 blocks. The current post-implementation results show that the full design still uses some DSP resources because wider fixed-point operations outside the native E2M1 multiply primitive can be inferred by HLS. Therefore, the correct claim is not that the entire implementation is DSP-free, but that the FP4 element multiply itself is LUT-based. Future revisions can further constrain adders and accumulators to fabric implementations using HLS binding directives and compare those resource results against the current implementation.

\subsection{Synthetic vs. Realistic Data Accuracy}
Testing in the present case is done through deterministic synthetic vectors. However, these vectors do not fully reflect the activation distributions of a complete language model. Distributions obtained from LLMs can contain high-magnitude outliers, widely observed in transformer-based models \cite{b2}, which may cause quantization error in regions where uniform or Gaussian synthetic distributions would not. This challenge is not unique to transformer architectures; recent linear and hybrid attention models, including GDN variants \cite{b3} and hybrid linear attention architectures \cite{b8}, similarly exhibit activation distributions that complicate low-bit quantization. This means that perplexity and downstream model quality have not yet been validated. Further testing must include real Qwen3-Next activation captures and perplexity evaluation on a canonical benchmark dataset.

\section{Future Work}
Future research should validate the proposed architecture using realistic Qwen3-Next activation captures to conduct full perplexity evaluations on standard benchmark datasets.
Additionally, the hardware implementation can be further optimized by enforcing stricter High-Level Synthesis (HLS) binding directives, ensuring that peripheral fixed-point operations are mapped to fabric logic rather than DSP blocks. Finally, deploying the accelerator onto physical FPGA hardware will allow for the measurement of true runtime overheads and provide empirical validation of the simulation-based power estimates.

\section{Conclusion}
This paper presents a simulation-backed native MXFP4 persistent-state datapath for one batch-one Gated DeltaNet decode layer on a Alveo U55C FPGA. The design keeps recurrent state on chip and uses block-scaled FP4 arithmetic to reduce decode latency and energy cost.
The default implementation achieves 51.472 microseconds per token in RTL cosimulation and 0.462 millijoules per token. Compared with the cited H100 baseline [1], [4] this corresponds to a 5.537x speedup and a 216.011x improvement in energy efficiency. The design also closes timing after Vivado implementation with 0.065 ns of positive slack while using 6.64% LUT, 3.67% BRAM, and 7.27% DSP. The current numerical results are limited to deterministic synthetic tests. Realistic Qwen3-Next activation capture and full-model perplexity evaluation remain future work. Therefore, the supported claim of this draft is a hardware-level result for RTL latency, implementation timing, resource use, power, and synthetic numerical fidelity, not a final Qwen3-Next model-quality claim.


\begin{thebibliography}{00}
\bibitem{b2} A. Vaswani et al., "Attention is all you need," in \textit{Advances in Neural Information Processing Systems (NeurIPS)}, vol. 30, 2017, pp. 5998--6008.
\bibitem{b6} H. Wu and K. Tu, "Layer-Condensed KV Cache for Efficient Inference of Large Language Models," arXiv preprint arXiv:2405.10637, 2024.
\bibitem{b9} Y. Zhang et al., "Kimi Linear: An Expressive, Efficient Attention Architecture," arXiv preprint arXiv:2510.26692, 2025.
\bibitem{b10} N. Gupta, P. Wang, R. Kannan, and V. K. Prasanna, "A Persistent-State Dataflow Accelerator for Memory-Bound Linear Attention Decode on FPGA," 2026.
\bibitem{b5} C. Wolters, X. Yang, U. Schlichtmann, and T. Suzumura, "Memory Is All You Need: An Overview of Compute-in-Memory Architectures for Accelerating Large Language Model Inference," arXiv preprint arXiv:2406.08413, 2024.
\bibitem{b12} B. D. Rouhani et al., "Microscaling data formats for deep learning," arXiv preprint arXiv:2310.10537, 2023.
\bibitem{b4} Y. Sheng et al., "High-throughput generative LLM inference with vLLM," in \textit{Proc. 29th ACM Symp. Operating Systems Principles (SOSP)}, 2023, pp. 163--181.
\bibitem{b1} R. Pope et al., "Efficiently scaling transformer inference," in \textit{Proc. Machine Learning and Systems (MLSys)}, vol. 5, 2023, pp. 606--624.
\bibitem{b7} P. G. Recasens et al., "Mind the Memory Gap: Unveiling GPU Bottlenecks in Large-Batch LLM Inference," in \textit{Proc. 2025 IEEE 18th Int. Conf. Cloud Computing (CLOUD)}, 2025, pp. 277--287.
\bibitem{b3} S. Yang, J. Kautz, and A. Hatamizadeh, "Gated Delta Networks: Improving Mamba2 with Delta Rule," arXiv preprint arXiv:2412.06464, 2024.
\bibitem{b13} V. Egiazarian et al., "Bridging the Gap Between Promise and Performance for Microscaling FP4 Quantization," arXiv preprint arXiv:2509.23202, 2025.
\bibitem{b14} M. Wipfli, G. Ä°slamoÄŸlu, N. K. Purayil, A. Garofalo, and L. Benini, "VMXDOTP: A RISC-V Vector ISA Extension for Efficient Microscaling (MX) Format Acceleration," arXiv preprint arXiv:2603.04979, 2026.
\bibitem{b15} M. Chen et al., "INT v.s. FP: A Comprehensive Study of Fine-Grained Low-bit Quantization Formats," arXiv preprint arXiv:2510.25602, 2025.
\bibitem{b8} D. Wang et al., "A Systematic Analysis of Hybrid Linear Attention," arXiv preprint arXiv:2507.06457, 2025.
\bibitem{b11} H. Guo et al., "Log-Linear Attention," arXiv preprint arXiv:2506.04761, 2025.
\end{thebibliography}

\end{document}
